% GENERATED FILE. Edit canonical lessons, then run npm run generate:books.
% canonical-source-hash: fnv1a64:5e5fa8a31f143eae
% canonical-lessons: ZH-C40-mai, ZH-W40-mai, ZH-C40-pingguo, ZH-C40-shuiguo, ZH-C40-mianbao, ZH-C40-jidan

\chapter{Fruit, Bread and Eggs}
\label{ch:zh-fruit-bread-and-eggs}
\hlchaptermodality{40}

\begin{chapteropening}
\textbf{By the end of this chapter:} \emph{I can say what I buy and name fruit, bread and eggs.}

Complete the last lesson of chapter 40: I can say what I buy and name fruit, bread and eggs.
\end{chapteropening}

\section[\texorpdfstring{\zh{买} (mǎi)}{买 (mǎi)}]{\zh{买} (mǎi) — to buy}

\label{lesson:ZH-C40-mai}



\begin{quote}
Before the new one: say the Mandarin for milk, then the Mandarin for tea.
\end{quote}

\subsection*{You'll want to know: \zh{买}}
\textbf{\zh{买}} — \emph{mǎi} — \textquotedblleft{}to buy\textquotedblright{}.

\begin{quote}
\textbf{\zh{你们买水吗}} \emph{nǐmen mǎi shuǐ ma} — Are you buying water?
\end{quote}

The new character, in pieces: \zh{买} — a lid with two dots under it, and below them a shape like \zh{大} with its last stroke cut down to a dot. The next lesson puts it in your hand.

\begin{cousinweb}
\textbf{\zh{买}} (\emph{mǎi}) is inherited from Middle Chinese \emph{meaX} (Guangyun, in Baxter's transcription), glossed 'to trade at market'.
\end{cousinweb}

\begin{grammarlens}[title={What you've built}]
One more word for everyday answers.
\end{grammarlens}

\subsection*{Your turn}
\begin{itemize}
\raggedright
  \item \emph{Say it:} \emph{mǎi}
  \item \emph{Say it:} \emph{mǎi}, once more
  \item \emph{Say it:} \emph{nǐmen mǎi shuǐ ma} — Are you buying water?
  \item \emph{Recall:} say the Mandarin for milk, then the Mandarin for tea, then say \emph{mǎi} again
\end{itemize}

\begin{tcolorbox}[breakable,colback=teal!4,colframe=teal!35!black,title={Before you move on}]
What does \zh{买} mean? (\textquotedblleft{}To buy\textquotedblright{}.) Say it once more.
\end{tcolorbox}
\section[\texorpdfstring{\zh{买} (mǎi)}{买 (mǎi)}]{\zh{买} — a lid, two dots, and a cut-down \zh{大}}

\label{lesson:ZH-W40-mai}



\begin{quote}
The word first, then the character inside it.

\begin{itemize}
\raggedright
  \item \emph{Recall:} say \emph{mǎi} — \textbf{R1}, one lesson back, before it decays
\end{itemize}

\begin{quote}
\zh{大} \zh{水}
\end{quote}

\emph{Read it:} each one aloud, with nothing said first — \emph{dà}, \emph{shuǐ}
\end{quote}

\subsection*{Script — a lid, two dots, and a cut-down \zh{大}}
\hlblockfigure{\detokenize{figures/ZH-W40-mai-filmstrip.pdf}}{How \zh{买} is written, stroke by stroke}

\begin{quote}
\zh{买}
\end{quote}

\textbf{Two pieces, top before bottom.}

\textbf{The top piece, three strokes.} A horizontal that hooks down at its right end, like a lid. Then two dots under it, one above the other.

\textbf{The bottom piece, three strokes.} A long horizontal, a left-falling stroke through it, and a dot on the right. It is \zh{大} with its last stroke cut down to a dot.

3 + 3 = six.

\subsection*{Your turn}
\begin{itemize}
\raggedright
  \item \emph{Write it:} \zh{买}, all six, with the model in front of you
  \item \emph{Cover:} write it again from memory, and count the pieces, not the strokes
\end{itemize}

\begin{tcolorbox}[breakable,colback=teal!4,colframe=teal!35!black,title={Before you move on}]
How many pieces in \zh{买}? (\textbf{Two} — the lid with its dots, then the cut-down \zh{大}.) What is under the lid? (\textbf{Two dots}.) What is the last stroke? (\textbf{A dot}, where \zh{大} has a long fall.)

Source: \href{https://raw.githubusercontent.com/chanind/hanzi-writer-data/68d10a4b21150cae5e1ebbd223eed289cf32d90c/data/\%E4\%B9\%B0.json}{Hanzi Writer Data \zh{买}, pinned stroke paths and medians}.
\end{tcolorbox}
\section[píngguǒ]{píngguǒ — apple}

\label{lesson:ZH-C40-pingguo}



\begin{quote}
Before the new one: say the Mandarin for tea, then the Mandarin for to buy.
\end{quote}

\subsection*{You'll want to know: píngguǒ}
\emph{píngguǒ} — \textquotedblleft{}apple\textquotedblright{}.

\begin{quote}
\emph{wǒ mǎi le sān ge píngguǒ} — I bought three apples.
\end{quote}

Its characters come later; for now it is a word for the ear.

\begin{cousinweb}
\emph{píngguǒ} is built from píng and guǒ 'fruit'. Guǒ is inherited from Middle Chinese \emph{kwaX} (Guangyun, in Baxter's transcription), glossed 'fruit of a tree'; the first character is listed there as \emph{bjaeng}, a name for several plants, duckweed among them, and how it came to name the apple is not settled.
\end{cousinweb}

\begin{grammarlens}[title={What you've built}]
One more word for everyday answers.
\end{grammarlens}

\subsection*{Your turn}
\begin{itemize}
\raggedright
  \item \emph{Say it:} \emph{píngguǒ}
  \item \emph{Say it:} \emph{píngguǒ}, once more
  \item \emph{Say it:} \emph{wǒ mǎi le sān ge píngguǒ} — I bought three apples.
  \item \emph{Recall:} say the Mandarin for tea, then the Mandarin for to buy, then say \emph{píngguǒ} again
\end{itemize}

\begin{tcolorbox}[breakable,colback=teal!4,colframe=teal!35!black,title={Before you move on}]
What does píngguǒ mean? (\textquotedblleft{}Apple\textquotedblright{}.) Say it once more.
\end{tcolorbox}
\section[shuǐguǒ]{shuǐguǒ — fruit}

\label{lesson:ZH-C40-shuiguo}



\begin{quote}
Before the new one: say the Mandarin for to buy, then the Mandarin for apple.
\end{quote}

\subsection*{You'll want to know: shuǐguǒ}
\emph{shuǐguǒ} — \textquotedblleft{}fruit\textquotedblright{}.

\begin{quote}
\emph{wǒ xǐhuan chī shuǐguǒ} — I like fruit.
\end{quote}

Its characters come later; for now it is a word for the ear.

Tone: \emph{shuǐguǒ} is two third tones in a row, so \emph{shuǐ} rises and only \emph{guǒ} keeps the dip. The spelling does not show the change; your voice makes it.

\begin{cousinweb}
\emph{shuǐguǒ} is built from \zh{水} 'water' and guǒ 'fruit', both inherited from Middle Chinese, \emph{sywijX} and \emph{kwaX} (Guangyun, in Baxter's transcription).
\end{cousinweb}

\begin{grammarlens}[title={What you've built}]
One more word for everyday answers.
\end{grammarlens}

\subsection*{Your turn}
\begin{itemize}
\raggedright
  \item \emph{Say it:} \emph{shuǐguǒ}
  \item \emph{Say it:} \emph{shuǐguǒ}, once more
  \item \emph{Say it:} \emph{wǒ xǐhuan chī shuǐguǒ} — I like fruit.
  \item \emph{Recall:} say the Mandarin for to buy, then the Mandarin for apple, then say \emph{shuǐguǒ} again
\end{itemize}

\begin{tcolorbox}[breakable,colback=teal!4,colframe=teal!35!black,title={Before you move on}]
What does shuǐguǒ mean? (\textquotedblleft{}Fruit\textquotedblright{}.) Say it once more.
\end{tcolorbox}
\section[miànbāo]{miànbāo — bread}

\label{lesson:ZH-C40-mianbao}



\begin{quote}
Before the new one: say the Mandarin for apple, then the Mandarin for fruit.
\end{quote}

\subsection*{You'll want to know: miànbāo}
\emph{miànbāo} — \textquotedblleft{}bread\textquotedblright{}.

\begin{quote}
\emph{wǒ chī le yí ge miànbāo} — I had a bread roll.
\end{quote}

Its characters come later; for now it is a word for the ear.

\begin{cousinweb}
\emph{miànbāo} is built from miàn 'flour' and bāo 'to wrap', both inherited from Middle Chinese, \emph{menH} and \emph{paew} (Guangyun, in Baxter's transcription).
\end{cousinweb}

\begin{grammarlens}[title={What you've built}]
One more word for everyday answers.
\end{grammarlens}

\subsection*{Your turn}
\begin{itemize}
\raggedright
  \item \emph{Say it:} \emph{miànbāo}
  \item \emph{Say it:} \emph{miànbāo}, once more
  \item \emph{Say it:} \emph{wǒ chī le yí ge miànbāo} — I had a bread roll.
  \item \emph{Recall:} say the Mandarin for apple, then the Mandarin for fruit, then say \emph{miànbāo} again
\end{itemize}

\begin{tcolorbox}[breakable,colback=teal!4,colframe=teal!35!black,title={Before you move on}]
What does miànbāo mean? (\textquotedblleft{}Bread\textquotedblright{}.) Say it once more.
\end{tcolorbox}
\section[jīdàn]{jīdàn — egg}

\label{lesson:ZH-C40-jidan}



\begin{quote}
Before the new one: say the Mandarin for fruit, then the Mandarin for bread.
\end{quote}

\subsection*{You'll want to know: jīdàn}
\emph{jīdàn} — \textquotedblleft{}egg\textquotedblright{}.

\begin{quote}
\emph{tā xǐhuan chī jīdàn} — He likes eggs.
\end{quote}

Its characters come later; for now it is a word for the ear.

\begin{cousinweb}
\emph{jīdàn} is built from jī 'chicken' and dàn 'egg'. The first is inherited from Middle Chinese \emph{kej} (Guangyun, in Baxter's transcription), which the dictionary calls the animal that knows the hour; no dictionary we use records an older form of the second, so its history is not settled.
\end{cousinweb}

\begin{grammarlens}[title={What you've built}]
One more word for everyday answers.
\end{grammarlens}

\subsection*{Your turn}
\begin{itemize}
\raggedright
  \item \emph{Say it:} \emph{jīdàn}
  \item \emph{Say it:} \emph{jīdàn}, once more
  \item \emph{Say it:} \emph{tā xǐhuan chī jīdàn} — He likes eggs.
  \item \emph{Recall:} say the Mandarin for fruit, then the Mandarin for bread, then say \emph{jīdàn} again
\end{itemize}

\begin{tcolorbox}[breakable,colback=teal!4,colframe=teal!35!black,title={Before you move on}]
What does jīdàn mean? (\textquotedblleft{}Egg\textquotedblright{}.) Say it once more.
\end{tcolorbox}
